% Part: computability
% Chapter: computability-theory
% Section: reducibility

\documentclass[../../../include/open-logic-section]{subfiles}

\begin{document}

\olfileid{cmp}{thy}{red}
\olsection{قابلية الاختزال}

\begin{explain}
نعرف الآن وجود مجموعة واحدة على الأقل، هي $K_0$، قابلة للتعداد حسابيًا
ولكنها غير قابلة للحساب. وينبغي أن يتضح وجود غيرها. وتوفر طريقة
الاختزال وسيلة قوية لبيان أن مجموعات أخرى لها هذه الخواص، من غير حاجة
إلى العودة كل مرة إلى المبادئ الأولى.

بوجه عام، «اختزال» مجموعة $A$ إلى مجموعة~$B$ هو طريقة لتحويل أجوبة
الأسئلة عن انتماء ال!!{element}s إلى~$B$ إلى أجوبة الأسئلة عن انتمائها
إلى~$A$. وسنركز على مفهوم يسمى «قابلية الاختزال متعددًا إلى واحد»، مع
أن ثمة مفاهيم أخرى كثيرة لقابلية الاختزال، ولها خواص متفاوتة. ومفاهيم
الاختزال مركزية أيضًا في دراسة التعقيد الحسابي، حيث يلزم النظر في مسائل
الكفاءة كذلك. فمثلًا يقال إن مجموعة «تامة في NP» إذا كانت في NP وأمكن
اختزال كل مسألة في NP إليها، باستعمال مفهوم اختزال يشبه المفهوم الموصوف
أدناه، مع الشرط الإضافي القائل بإمكان حساب الاختزال في زمن كثير الحدود.

وقد استعملنا مفهوم الاختزال ضمنًا من قبل. عرف المجموعة~$K$ بـ
\[
K = \Setabs{x}{\cfind{x}(x) \fdefined},
\]
أي $K=\Setabs{x}{x\in W_x}$. ويبين برهاننا على عدم قابلية مشكلة
التوقف للحل (\olref[hlt]{thm:halting-problem})، على نحو مباشر جدًا، أن
$K$ غير قابلة للحساب. وتذكر أن $K_0$ هي المجموعة
\[
K_0 = \Setabs{\tuple{e, x}}{\cfind{e}(x) \fdefined },
\]
أي $K_0=\Setabs{\tuple{e,x}}{x\in W_e}$. ويسهل توسيع أي برهان على
عدم قابلية~$K$ للحساب إلى برهان على عدم قابلية~$K_0$ للحساب: فلو كانت
$K_0$ قابلة للحساب، لأمكننا تقرير ما إذا كان !!a{element}~$x$ في $K$
بمجرد السؤال عما إذا كان الزوج $\tuple{x,x}$ في $K_0$. والدالة~$f$
التي تعين $x$ إلى $\tuple{x,x}$ مثال \emph{لاختزال} $K$ إلى $K_0$.
\end{explain}

\begin{defn}
لتكن $A$ و$B$ مجموعتين من الأعداد الطبيعية. تكون الدالة القابلة للحساب
~$f\colon\Nat\to\Nat$ \emph{اختزالًا متعددًا إلى واحد} لـ$A$ إلى~$B$
إذا وفقط إذا كان، لكل عدد طبيعي~$x$،
\[
x \in A \quad \text{إذا وفقط إذا} \quad f(x) \in B.
\]
وإذا وجد اختزال $f$ كهذا، قلنا إن $A$ \emph{قابلة للاختزال متعددًا
إلى واحد} إلى~$B$، وكتبنا $A\leq_m B$. وإذا كانت $A$ قابلة للاختزال
متعددًا إلى واحد إلى $B$ والعكس، قيل إن $A$ و$B$ \emph{متكافئتان
متعددًا إلى واحد}، وكتبنا $A\equiv_m B$.
\end{defn}

إذا صادف أن كانت الدالة $f$ في التعريف أعلاه أحادية، قيل إن $A$
\emph{قابلة للاختزال واحدًا إلى واحد} إلى~$B$. وتحقق معظم الاختزالات
الموصوفة أدناه هذا الشرط الأقوى، لكننا لن نستعمل هذه الحقيقة.

\begin{digress}
صحيح، وإن لم يكن واضحًا البتة، أن قابلية الاختزال واحدًا إلى واحد شرط
أقوى فعلًا من قابلية الاختزال متعددًا إلى واحد. وبعبارة أخرى، توجد
مجموعتان لا نهائيتان $A$ و~$B$ تكون $A$ فيهما قابلة للاختزال متعددًا
إلى واحد إلى~$B$، ولكنها ليست قابلة للاختزال واحدًا إلى واحد إلى~$B$.
\end{digress}

\end{document}
